% ─────────────────────────────────────────────────────────────────────────────
% K04 §4.2 Bouncing Ball Downstairs — QSS1 Experiment Results
% Included by docs/main.tex via % ─────────────────────────────────────────────────────────────────────────────
% K04 §4.2 Bouncing Ball Downstairs — QSS1 Experiment Results
% Included by docs/main.tex via % ─────────────────────────────────────────────────────────────────────────────
% K04 §4.2 Bouncing Ball Downstairs — QSS1 Experiment Results
% Included by docs/main.tex via % ─────────────────────────────────────────────────────────────────────────────
% K04 §4.2 Bouncing Ball Downstairs — QSS1 Experiment Results
% Included by docs/main.tex via \input{../qss1-ball-downstairs/results}
% ─────────────────────────────────────────────────────────────────────────────

\subsection*{Overview}

This experiment replicates Example~4.2 of Kofman~\cite{K04}: a ball
moves in the vertical plane and bounces down a staircase whose steps
are $1\,\text{m}$ wide and $1\,\text{m}$ tall.  Contact with each
step is modelled by a spring--damper law rather than an instantaneous
velocity reversal, so the system is a hybrid ODE with a continuously
varying stiffness switch.

The simulation uses a QSS1 (first-order Quantized State System)
integrator for each state variable, coupled via Classic DEVS internal
couplings in \texttt{cadmium}.

% ─────────────────────────────────────────────────────────────────────────────
\subsection*{ODEs}

\[
  \dot{x}  = v_x,\qquad
  \dot{v}_x = -\tfrac{b_a}{m}\,v_x
\]
\[
  \dot{y}  = v_y,\qquad
  \dot{v}_y = -g
    - \tfrac{b_a}{m}\,v_y
    - s_w\!\left[\tfrac{b}{m}\,v_y + \tfrac{k}{m}\,(y - y_{\rm floor})\right]
\]

where $y_{\rm floor}(x) = \lfloor h + 1 - x \rfloor$ is the height of
the stair below position $x$, and

\[
  s_w =
  \begin{cases}
    1 & y \le y_{\rm floor}(x) \quad\text{(contact)}, \\
    0 & \text{otherwise}.
  \end{cases}
\]

% ─────────────────────────────────────────────────────────────────────────────
\subsection*{Parameters and Initial Conditions}

\begin{table}[h!]
\centering
\caption{Physical parameters (K04 Table~1).}
\label{tab:k04-params}
\begin{tabular}{clll}
\toprule
Symbol & Value & Units & Description \\
\midrule
$m$   & $1$        & kg           & ball mass \\
$k$   & $100\,000$ & N/m          & contact spring stiffness \\
$b$   & $30$       & N\,s/m       & contact damping coefficient \\
$b_a$ & $0.1$      & N\,s/m       & aerodynamic drag coefficient \\
$g$   & $9.8$      & m/s$^2$      & gravitational acceleration \\
$h$   & $10$       & m            & height of the first stair step \\
\bottomrule
\end{tabular}
\end{table}

\begin{table}[h!]
\centering
\caption{Initial conditions (K04 p.~19).}
\label{tab:k04-ic}
\begin{tabular}{cll}
\toprule
Variable & Value & Units \\
\midrule
$x(0)$   & $0.575$ & m     \\
$v_x(0)$ & $0.5$   & m/s   \\
$y(0)$   & $10.5$  & m     \\
$v_y(0)$ & $0$     & m/s   \\
\bottomrule
\end{tabular}
\end{table}

\begin{table}[h!]
\centering
\caption{QSS1 quantization parameters.}
\label{tab:k04-quant}
\begin{tabular}{clll}
\toprule
Variable & $\Delta q$ & Units & Note \\
\midrule
$x$   & $0.001$ & m   & \\
$v_x$ & $0.01$  & m/s & \\
$y$   & $0.001$ & m   & K04 uses $10^{-4}$; increased to suppress Zeno
                         micro-bouncing ($\delta \approx 8\times10^{-4}$\,m) \\
$v_y$ & $0.01$  & m/s & \\
\bottomrule
\end{tabular}
\end{table}

% ─────────────────────────────────────────────────────────────────────────────
\subsection*{Simulation Results}

Figure~\ref{fig:k04-trajectory} shows the QSS1 trajectories of the
horizontal position $x(t)$ and vertical position $y(t)$ over
$t\in[0,10]\,\text{s}$, matching the timeframe of K04 Figures~11--12.
The ball descends the staircase, losing height at each bounce while
advancing horizontally with exponentially damped velocity.
Dashed horizontal lines in the lower panel mark the floor levels of
the visited stair steps ($y = 7, 8, 9, 10$\,m).

Trajectory data is generated by \texttt{extract\_trajectories.py},
which streams NDJSON from the binary and writes
\texttt{figures/traj\_x.dat} and \texttt{figures/traj\_y.dat}.
The $y$ series is thinned to one sample per $2\,\text{ms}$ of
simulation time (4\,796 points) for compilation speed; the $x$ series
retains all 3\,226 quantization events.

\begin{figure}[h!]
\centering
\begin{tikzpicture}
\begin{groupplot}[
    group style={
        group size=1 by 2,
        xlabels at=edge bottom,
        xticklabels at=edge bottom,
        vertical sep=14pt,
    },
    width=\linewidth,
    height=5.2cm,
    xmin=0, xmax=10,
    grid=both,
    grid style={line width=0.3pt, draw=gray!30},
    major grid style={line width=0.4pt, draw=gray!50},
    tick align=outside,
    tick pos=left,
    xlabel={$t$ (s)},
    every axis label/.style={font=\small},
    every tick label/.style={font=\small},
]
% ── x(t): horizontal position ────────────────────────────────────────────────
\nextgroupplot[
    ylabel={$x$ (m)},
    ymin=0, ymax=5,
]
\addplot[blue!70!black, line width=0.6pt]
    table[x=t, y=q]{../qss1-ball-downstairs/figures/traj_x.dat};

% ── y(t): vertical position with stair-floor reference lines ─────────────────
\nextgroupplot[
    ylabel={$y$ (m)},
    ymin=6, ymax=11,
]
\addplot[blue!70!black, line width=0.6pt]
    table[x=t, y=q]{../qss1-ball-downstairs/figures/traj_y.dat};
% Stair floor levels visited during t=[0,10]
\foreach \lvl in {7,8,9,10} {
    \addplot[gray, dashed, line width=0.4pt]
        coordinates {(0,\lvl)(10,\lvl)};
}
\end{groupplot}
\end{tikzpicture}
\caption{QSS1 trajectory for K04~§4.2 over $t\in[0,10]\,\text{s}$.
  \emph{Top:} horizontal position $x(t)$.
  \emph{Bottom:} vertical position $y(t)$; dashed lines show stair-step
  floor levels at $y=7,8,9,10$\,m.}
\label{fig:k04-trajectory}
\end{figure}

% ─────────────────────────────────────────────────────────────────────────────

\subsection*{Overview}

This experiment replicates Example~4.2 of Kofman~\cite{K04}: a ball
moves in the vertical plane and bounces down a staircase whose steps
are $1\,\text{m}$ wide and $1\,\text{m}$ tall.  Contact with each
step is modelled by a spring--damper law rather than an instantaneous
velocity reversal, so the system is a hybrid ODE with a continuously
varying stiffness switch.

The simulation uses a QSS1 (first-order Quantized State System)
integrator for each state variable, coupled via Classic DEVS internal
couplings in \texttt{cadmium}.

% ─────────────────────────────────────────────────────────────────────────────
\subsection*{ODEs}

\[
  \dot{x}  = v_x,\qquad
  \dot{v}_x = -\tfrac{b_a}{m}\,v_x
\]
\[
  \dot{y}  = v_y,\qquad
  \dot{v}_y = -g
    - \tfrac{b_a}{m}\,v_y
    - s_w\!\left[\tfrac{b}{m}\,v_y + \tfrac{k}{m}\,(y - y_{\rm floor})\right]
\]

where $y_{\rm floor}(x) = \lfloor h + 1 - x \rfloor$ is the height of
the stair below position $x$, and

\[
  s_w =
  \begin{cases}
    1 & y \le y_{\rm floor}(x) \quad\text{(contact)}, \\
    0 & \text{otherwise}.
  \end{cases}
\]

% ─────────────────────────────────────────────────────────────────────────────
\subsection*{Parameters and Initial Conditions}

\begin{table}[h!]
\centering
\caption{Physical parameters (K04 Table~1).}
\label{tab:k04-params}
\begin{tabular}{clll}
\toprule
Symbol & Value & Units & Description \\
\midrule
$m$   & $1$        & kg           & ball mass \\
$k$   & $100\,000$ & N/m          & contact spring stiffness \\
$b$   & $30$       & N\,s/m       & contact damping coefficient \\
$b_a$ & $0.1$      & N\,s/m       & aerodynamic drag coefficient \\
$g$   & $9.8$      & m/s$^2$      & gravitational acceleration \\
$h$   & $10$       & m            & height of the first stair step \\
\bottomrule
\end{tabular}
\end{table}

\begin{table}[h!]
\centering
\caption{Initial conditions (K04 p.~19).}
\label{tab:k04-ic}
\begin{tabular}{cll}
\toprule
Variable & Value & Units \\
\midrule
$x(0)$   & $0.575$ & m     \\
$v_x(0)$ & $0.5$   & m/s   \\
$y(0)$   & $10.5$  & m     \\
$v_y(0)$ & $0$     & m/s   \\
\bottomrule
\end{tabular}
\end{table}

\begin{table}[h!]
\centering
\caption{QSS1 quantization parameters.}
\label{tab:k04-quant}
\begin{tabular}{clll}
\toprule
Variable & $\Delta q$ & Units & Note \\
\midrule
$x$   & $0.001$ & m   & \\
$v_x$ & $0.01$  & m/s & \\
$y$   & $0.001$ & m   & K04 uses $10^{-4}$; increased to suppress Zeno
                         micro-bouncing ($\delta \approx 8\times10^{-4}$\,m) \\
$v_y$ & $0.01$  & m/s & \\
\bottomrule
\end{tabular}
\end{table}

% ─────────────────────────────────────────────────────────────────────────────
\subsection*{Simulation Results}

Figure~\ref{fig:k04-trajectory} shows the QSS1 trajectories of the
horizontal position $x(t)$ and vertical position $y(t)$ over
$t\in[0,10]\,\text{s}$, matching the timeframe of K04 Figures~11--12.
The ball descends the staircase, losing height at each bounce while
advancing horizontally with exponentially damped velocity.
Dashed horizontal lines in the lower panel mark the floor levels of
the visited stair steps ($y = 7, 8, 9, 10$\,m).

Trajectory data is generated by \texttt{extract\_trajectories.py},
which streams NDJSON from the binary and writes
\texttt{figures/traj\_x.dat} and \texttt{figures/traj\_y.dat}.
The $y$ series is thinned to one sample per $2\,\text{ms}$ of
simulation time (4\,796 points) for compilation speed; the $x$ series
retains all 3\,226 quantization events.

\begin{figure}[h!]
\centering
\begin{tikzpicture}
\begin{groupplot}[
    group style={
        group size=1 by 2,
        xlabels at=edge bottom,
        xticklabels at=edge bottom,
        vertical sep=14pt,
    },
    width=\linewidth,
    height=5.2cm,
    xmin=0, xmax=10,
    grid=both,
    grid style={line width=0.3pt, draw=gray!30},
    major grid style={line width=0.4pt, draw=gray!50},
    tick align=outside,
    tick pos=left,
    xlabel={$t$ (s)},
    every axis label/.style={font=\small},
    every tick label/.style={font=\small},
]
% ── x(t): horizontal position ────────────────────────────────────────────────
\nextgroupplot[
    ylabel={$x$ (m)},
    ymin=0, ymax=5,
]
\addplot[blue!70!black, line width=0.6pt]
    table[x=t, y=q]{../qss1-ball-downstairs/figures/traj_x.dat};

% ── y(t): vertical position with stair-floor reference lines ─────────────────
\nextgroupplot[
    ylabel={$y$ (m)},
    ymin=6, ymax=11,
]
\addplot[blue!70!black, line width=0.6pt]
    table[x=t, y=q]{../qss1-ball-downstairs/figures/traj_y.dat};
% Stair floor levels visited during t=[0,10]
\foreach \lvl in {7,8,9,10} {
    \addplot[gray, dashed, line width=0.4pt]
        coordinates {(0,\lvl)(10,\lvl)};
}
\end{groupplot}
\end{tikzpicture}
\caption{QSS1 trajectory for K04~§4.2 over $t\in[0,10]\,\text{s}$.
  \emph{Top:} horizontal position $x(t)$.
  \emph{Bottom:} vertical position $y(t)$; dashed lines show stair-step
  floor levels at $y=7,8,9,10$\,m.}
\label{fig:k04-trajectory}
\end{figure}

% ─────────────────────────────────────────────────────────────────────────────

\subsection*{Overview}

This experiment replicates Example~4.2 of Kofman~\cite{K04}: a ball
moves in the vertical plane and bounces down a staircase whose steps
are $1\,\text{m}$ wide and $1\,\text{m}$ tall.  Contact with each
step is modelled by a spring--damper law rather than an instantaneous
velocity reversal, so the system is a hybrid ODE with a continuously
varying stiffness switch.

The simulation uses a QSS1 (first-order Quantized State System)
integrator for each state variable, coupled via Classic DEVS internal
couplings in \texttt{cadmium}.

% ─────────────────────────────────────────────────────────────────────────────
\subsection*{ODEs}

\[
  \dot{x}  = v_x,\qquad
  \dot{v}_x = -\tfrac{b_a}{m}\,v_x
\]
\[
  \dot{y}  = v_y,\qquad
  \dot{v}_y = -g
    - \tfrac{b_a}{m}\,v_y
    - s_w\!\left[\tfrac{b}{m}\,v_y + \tfrac{k}{m}\,(y - y_{\rm floor})\right]
\]

where $y_{\rm floor}(x) = \lfloor h + 1 - x \rfloor$ is the height of
the stair below position $x$, and

\[
  s_w =
  \begin{cases}
    1 & y \le y_{\rm floor}(x) \quad\text{(contact)}, \\
    0 & \text{otherwise}.
  \end{cases}
\]

% ─────────────────────────────────────────────────────────────────────────────
\subsection*{Parameters and Initial Conditions}

\begin{table}[h!]
\centering
\caption{Physical parameters (K04 Table~1).}
\label{tab:k04-params}
\begin{tabular}{clll}
\toprule
Symbol & Value & Units & Description \\
\midrule
$m$   & $1$        & kg           & ball mass \\
$k$   & $100\,000$ & N/m          & contact spring stiffness \\
$b$   & $30$       & N\,s/m       & contact damping coefficient \\
$b_a$ & $0.1$      & N\,s/m       & aerodynamic drag coefficient \\
$g$   & $9.8$      & m/s$^2$      & gravitational acceleration \\
$h$   & $10$       & m            & height of the first stair step \\
\bottomrule
\end{tabular}
\end{table}

\begin{table}[h!]
\centering
\caption{Initial conditions (K04 p.~19).}
\label{tab:k04-ic}
\begin{tabular}{cll}
\toprule
Variable & Value & Units \\
\midrule
$x(0)$   & $0.575$ & m     \\
$v_x(0)$ & $0.5$   & m/s   \\
$y(0)$   & $10.5$  & m     \\
$v_y(0)$ & $0$     & m/s   \\
\bottomrule
\end{tabular}
\end{table}

\begin{table}[h!]
\centering
\caption{QSS1 quantization parameters.}
\label{tab:k04-quant}
\begin{tabular}{clll}
\toprule
Variable & $\Delta q$ & Units & Note \\
\midrule
$x$   & $0.001$ & m   & \\
$v_x$ & $0.01$  & m/s & \\
$y$   & $0.001$ & m   & K04 uses $10^{-4}$; increased to suppress Zeno
                         micro-bouncing ($\delta \approx 8\times10^{-4}$\,m) \\
$v_y$ & $0.01$  & m/s & \\
\bottomrule
\end{tabular}
\end{table}

% ─────────────────────────────────────────────────────────────────────────────
\subsection*{Simulation Results}

Figure~\ref{fig:k04-trajectory} shows the QSS1 trajectories of the
horizontal position $x(t)$ and vertical position $y(t)$ over
$t\in[0,10]\,\text{s}$, matching the timeframe of K04 Figures~11--12.
The ball descends the staircase, losing height at each bounce while
advancing horizontally with exponentially damped velocity.
Dashed horizontal lines in the lower panel mark the floor levels of
the visited stair steps ($y = 7, 8, 9, 10$\,m).

Trajectory data is generated by \texttt{extract\_trajectories.py},
which streams NDJSON from the binary and writes
\texttt{figures/traj\_x.dat} and \texttt{figures/traj\_y.dat}.
The $y$ series is thinned to one sample per $2\,\text{ms}$ of
simulation time (4\,796 points) for compilation speed; the $x$ series
retains all 3\,226 quantization events.

\begin{figure}[h!]
\centering
\begin{tikzpicture}
\begin{groupplot}[
    group style={
        group size=1 by 2,
        xlabels at=edge bottom,
        xticklabels at=edge bottom,
        vertical sep=14pt,
    },
    width=\linewidth,
    height=5.2cm,
    xmin=0, xmax=10,
    grid=both,
    grid style={line width=0.3pt, draw=gray!30},
    major grid style={line width=0.4pt, draw=gray!50},
    tick align=outside,
    tick pos=left,
    xlabel={$t$ (s)},
    every axis label/.style={font=\small},
    every tick label/.style={font=\small},
]
% ── x(t): horizontal position ────────────────────────────────────────────────
\nextgroupplot[
    ylabel={$x$ (m)},
    ymin=0, ymax=5,
]
\addplot[blue!70!black, line width=0.6pt]
    table[x=t, y=q]{../qss1-ball-downstairs/figures/traj_x.dat};

% ── y(t): vertical position with stair-floor reference lines ─────────────────
\nextgroupplot[
    ylabel={$y$ (m)},
    ymin=6, ymax=11,
]
\addplot[blue!70!black, line width=0.6pt]
    table[x=t, y=q]{../qss1-ball-downstairs/figures/traj_y.dat};
% Stair floor levels visited during t=[0,10]
\foreach \lvl in {7,8,9,10} {
    \addplot[gray, dashed, line width=0.4pt]
        coordinates {(0,\lvl)(10,\lvl)};
}
\end{groupplot}
\end{tikzpicture}
\caption{QSS1 trajectory for K04~§4.2 over $t\in[0,10]\,\text{s}$.
  \emph{Top:} horizontal position $x(t)$.
  \emph{Bottom:} vertical position $y(t)$; dashed lines show stair-step
  floor levels at $y=7,8,9,10$\,m.}
\label{fig:k04-trajectory}
\end{figure}

% ─────────────────────────────────────────────────────────────────────────────

\subsection*{Overview}

This experiment replicates Example~4.2 of Kofman~\cite{K04}: a ball
moves in the vertical plane and bounces down a staircase whose steps
are $1\,\text{m}$ wide and $1\,\text{m}$ tall.  Contact with each
step is modelled by a spring--damper law rather than an instantaneous
velocity reversal, so the system is a hybrid ODE with a continuously
varying stiffness switch.

The simulation uses a QSS1 (first-order Quantized State System)
integrator for each state variable, coupled via Classic DEVS internal
couplings in \texttt{cadmium}.

% ─────────────────────────────────────────────────────────────────────────────
\subsection*{ODEs}

\[
  \dot{x}  = v_x,\qquad
  \dot{v}_x = -\tfrac{b_a}{m}\,v_x
\]
\[
  \dot{y}  = v_y,\qquad
  \dot{v}_y = -g
    - \tfrac{b_a}{m}\,v_y
    - s_w\!\left[\tfrac{b}{m}\,v_y + \tfrac{k}{m}\,(y - y_{\rm floor})\right]
\]

where $y_{\rm floor}(x) = \lfloor h + 1 - x \rfloor$ is the height of
the stair below position $x$, and

\[
  s_w =
  \begin{cases}
    1 & y \le y_{\rm floor}(x) \quad\text{(contact)}, \\
    0 & \text{otherwise}.
  \end{cases}
\]

% ─────────────────────────────────────────────────────────────────────────────
\subsection*{Parameters and Initial Conditions}

\begin{table}[h!]
\centering
\caption{Physical parameters (K04 Table~1).}
\label{tab:k04-params}
\begin{tabular}{clll}
\toprule
Symbol & Value & Units & Description \\
\midrule
$m$   & $1$        & kg           & ball mass \\
$k$   & $100\,000$ & N/m          & contact spring stiffness \\
$b$   & $30$       & N\,s/m       & contact damping coefficient \\
$b_a$ & $0.1$      & N\,s/m       & aerodynamic drag coefficient \\
$g$   & $9.8$      & m/s$^2$      & gravitational acceleration \\
$h$   & $10$       & m            & height of the first stair step \\
\bottomrule
\end{tabular}
\end{table}

\begin{table}[h!]
\centering
\caption{Initial conditions (K04 p.~19).}
\label{tab:k04-ic}
\begin{tabular}{cll}
\toprule
Variable & Value & Units \\
\midrule
$x(0)$   & $0.575$ & m     \\
$v_x(0)$ & $0.5$   & m/s   \\
$y(0)$   & $10.5$  & m     \\
$v_y(0)$ & $0$     & m/s   \\
\bottomrule
\end{tabular}
\end{table}

\begin{table}[h!]
\centering
\caption{QSS1 quantization parameters.}
\label{tab:k04-quant}
\begin{tabular}{clll}
\toprule
Variable & $\Delta q$ & Units & Note \\
\midrule
$x$   & $0.001$ & m   & \\
$v_x$ & $0.01$  & m/s & \\
$y$   & $0.001$ & m   & K04 uses $10^{-4}$; increased to suppress Zeno
                         micro-bouncing ($\delta \approx 8\times10^{-4}$\,m) \\
$v_y$ & $0.01$  & m/s & \\
\bottomrule
\end{tabular}
\end{table}

% ─────────────────────────────────────────────────────────────────────────────
\subsection*{Simulation Results}

Figure~\ref{fig:k04-trajectory} shows the QSS1 trajectories of the
horizontal position $x(t)$ and vertical position $y(t)$ over
$t\in[0,10]\,\text{s}$, matching the timeframe of K04 Figures~11--12.
The ball descends the staircase, losing height at each bounce while
advancing horizontally with exponentially damped velocity.
Dashed horizontal lines in the lower panel mark the floor levels of
the visited stair steps ($y = 7, 8, 9, 10$\,m).

Trajectory data is generated by \texttt{extract\_trajectories.py},
which streams NDJSON from the binary and writes
\texttt{figures/traj\_x.dat} and \texttt{figures/traj\_y.dat}.
The $y$ series is thinned to one sample per $2\,\text{ms}$ of
simulation time (4\,796 points) for compilation speed; the $x$ series
retains all 3\,226 quantization events.

\begin{figure}[h!]
\centering
\begin{tikzpicture}
\begin{groupplot}[
    group style={
        group size=1 by 2,
        xlabels at=edge bottom,
        xticklabels at=edge bottom,
        vertical sep=14pt,
    },
    width=\linewidth,
    height=5.2cm,
    xmin=0, xmax=10,
    grid=both,
    grid style={line width=0.3pt, draw=gray!30},
    major grid style={line width=0.4pt, draw=gray!50},
    tick align=outside,
    tick pos=left,
    xlabel={$t$ (s)},
    every axis label/.style={font=\small},
    every tick label/.style={font=\small},
]
% ── x(t): horizontal position ────────────────────────────────────────────────
\nextgroupplot[
    ylabel={$x$ (m)},
    ymin=0, ymax=5,
]
\addplot[blue!70!black, line width=0.6pt]
    table[x=t, y=q]{../qss1-ball-downstairs/figures/traj_x.dat};

% ── y(t): vertical position with stair-floor reference lines ─────────────────
\nextgroupplot[
    ylabel={$y$ (m)},
    ymin=6, ymax=11,
]
\addplot[blue!70!black, line width=0.6pt]
    table[x=t, y=q]{../qss1-ball-downstairs/figures/traj_y.dat};
% Stair floor levels visited during t=[0,10]
\foreach \lvl in {7,8,9,10} {
    \addplot[gray, dashed, line width=0.4pt]
        coordinates {(0,\lvl)(10,\lvl)};
}
\end{groupplot}
\end{tikzpicture}
\caption{QSS1 trajectory for K04~§4.2 over $t\in[0,10]\,\text{s}$.
  \emph{Top:} horizontal position $x(t)$.
  \emph{Bottom:} vertical position $y(t)$; dashed lines show stair-step
  floor levels at $y=7,8,9,10$\,m.}
\label{fig:k04-trajectory}
\end{figure}
